\chapter{研究背景と電力市場の構造変化}

本論文は、日本の新電力企業の倒産リスク構造を、Merton 型倒産距離モデルおよび
Ordered Logit モデルを用いて分析することを目的とする。本章では、その前提となる
世界および日本の電力市場の構造変化、新電力のビジネスモデル、倒産・退出の定義、
リスク構造について整理し、第2章の問題意識および研究目的へと接続する。

% ============================================================
% 1.1 世界のエネルギー情勢と電力市場の変容（拡張版）
% ============================================================

\section{世界のエネルギー情勢と電力市場の変容}

本節では、世界的なエネルギー転換の潮流と、それに伴う電力市場の構造変化を詳細に整理する。
既存の議論では、再エネ拡大・地政学リスク・燃料価格変動・需給逼迫・市場自由化などが
断片的に扱われることが多いが、本研究では新電力の倒産リスク分析に直結するため、
これらの要因を体系的に統合し、電力市場の「構造的変容」を明確に位置づける。

世界の電力市場は、以下の三つの外生ショックによって構造的に変化している。

\begin{enumerate}
  \item 供給側ショック：再エネ比率の急上昇と変動性の増大  
  \item 需要側ショック：AI・データセンターによる電力需要の急増  
  \item 地政学ショック：燃料価格の急騰と国際情勢の不確実性  
\end{enumerate}

これら三つのショックは、従来の集中型電力システムでは吸収しきれない規模に達しており、
自由化市場における価格形成を不安定化させ、小売事業者の経営リスクを増幅させている。
特に新電力はアセットライト型であり、外部ショックの影響を直接受けやすい構造を持つため、
世界的潮流の理解は日本市場の分析に不可欠である。

以下では、供給側・地政学・需要側の三つの観点から、世界の構造変化を詳細に整理する。

% ------------------------------------------------------------
% 1.1.1 供給側の構造変化：再エネ拡大と電力システムの変動性
% ------------------------------------------------------------

\subsection{供給側の構造変化：再エネ拡大と電力システムの変動性}

21世紀以降、世界の電力供給構造は、化石燃料依存型から脱炭素型へと急速に転換している。
特に太陽光・風力といった変動性再エネ（VRE: Variable Renewable Energy）の導入は、
電力システムの需給調整を根本的に変化させた。

IEA の統計によれば、2023 年の世界の再エネ発電容量は 510 GW 増加し、
過去最大の伸びを記録した\parencite{IEA2023Renewables}。
太陽光は年間 380 GW 以上増加し、世界の新規電源の 75\% を占める。

欧州では、再エネ比率が 2023 年に \textbf{44\%} に達し、
そのうち太陽光と風力が 25\% を占めるなど、電力供給の中心が急速に再エネへ移行している。
（※あなたの指摘箇所を修復済み）

再エネ拡大は脱炭素化の推進力である一方、以下の構造的課題を生む。

\begin{itemize}
  \item 出力変動が大きく、需給調整力（火力・蓄電池・DR）が不足しやすい  
  \item 天候依存性が高く、短期的な価格変動を増幅する  
  \item 系統制約（カーブアウト・出力抑制）が頻発し、市場価格形成を歪める  
\end{itemize}

欧州では、風力発電の出力が 1 日で 200 GWh 以上変動する事例が報告されており、
需給調整の難易度は従来の火力中心のシステムとは比較にならない。

OECD の推計では、2030 年までに欧州で 80 GW の調整力が不足する可能性があり、
再エネ拡大は電力システムの変動性と不確実性を構造的に増幅させている\parencite{OECD2022Security}。

この供給側ショックは、自由化市場における価格変動性を高め、
小売事業者の調達リスクを増幅させる。

% ------------------------------------------------------------
% 1.1.2 政治・地政学リスク：燃料価格と国際情勢の不確実性
% ------------------------------------------------------------

\subsection{政治・地政学リスク：燃料価格と国際情勢の不確実性}

世界の電力市場を不安定化させている第二の要因が、政治・地政学リスクである。
中東情勢の緊張、ロシア・ウクライナ戦争、米中対立などは、
化石燃料の供給網に直接的な影響を与え、燃料価格を急騰させた。

2021〜2022 年の燃料価格高騰では、欧州の LNG 価格（TTF 指標）が  
\textbf{20 €/MWh → 200 €/MWh 超}へと 10 倍以上に急騰し、  
石炭価格は 3 倍、原油価格も 2 倍に上昇した\parencite{IEA2022GasMarket}。

ロシア産ガス依存度は欧州で 40\% から 10\% 以下へ急減し、
その代替調達コストが電力価格の上昇圧力となった。

IEA によれば、2021〜2023 年の間に世界各国で延べ 70 回以上の需給逼迫警報が発令され、
地政学ショックが電力市場の不安定化を加速させている\parencite{IEA2024Electricity}。

このショックは、燃料調達力の弱い小売事業者にとって致命的であり、
特に新電力はヘッジ手段が限られるため、価格変動の影響を直接受けやすい。

% ------------------------------------------------------------
% 1.1.3 需要側の構造変化：AI・データセンターによる電力需要の急増
% ------------------------------------------------------------

\subsection{需要側の構造変化：AI・データセンターによる電力需要の急増}

近年では、供給側の変動性に加えて、需要側の急増が新たな構造的課題となっている。
特に AI の急速な普及に伴い、データセンターが世界の電力需要を押し上げる主要因となっている。

IEA の最新報告によれば、世界のデータセンターの電力消費は  
2024 年の 460 TWh から、2030 年には 1,000 TWh を超え、  
2035 年には 1,300 TWh に達する見通しである\parencite{IEA2025EnergySupplyAI}。

これは、日本全体の年間電力消費量に匹敵する規模であり、
電力システムに対する負荷は極めて大きい。

さらに、AI 最適化型データセンターに限れば、2030 年までに電力需要が 4 倍以上に増加するとの推計もある\parencite{IEA2025AISurge}。

米国では、2030 年までの電力需要増加の約半分がデータセンター由来と見込まれており、
AI の普及が電力需要構造を根本的に変えつつある。

この需要側ショックは、需給逼迫の頻度を高め、スポット価格の変動性を増幅し、
小売事業者の調達リスクをさらに大きくする。

% ------------------------------------------------------------
% 1.1.4 電力市場自由化の国際的潮流（拡張版）
% ------------------------------------------------------------

\subsection{電力市場自由化の国際的潮流}

世界の電力市場自由化は、各国の制度設計やエネルギー政策の違いによって多様な形態をとる。
本節では、世界の自由化モデルを「欧州モデル」「北米モデル」「アジアモデル」「新興国モデル」の
四つに類型化し、それぞれの特徴を整理する。

% --- 欧州モデル ---
\subsubsection{（1）欧州モデル：完全自由化と市場統合}

欧州連合（EU）は、世界で最も進んだ電力市場自由化を実現している。
発電・送電・小売の完全分離（アンバンドリング）が徹底され、
卸電力市場はスポット・先渡・時間前市場が高度に整備されている。

欧州の再エネ比率は 2023 年に \textbf{44\%} に達し、
変動性電源の増加に伴い市場価格の変動幅も拡大している。

欧州の小売スイッチング率は年間 \textbf{10〜20\%} と高く、
競争が活発である一方、価格変動リスクが小売事業者に直接転嫁される構造が形成されている。
（※あなたの指摘箇所を修復済み）

Nord Pool では、日間の価格変動が ±200 €/MWh に達する事例も報告されており、
需給調整力の確保が重要な政策課題となっている。

% --- 北米モデル ---
\subsubsection{（2）北米モデル：部分自由化と州別市場設計}

北米では、州ごとに自由化の度合いが大きく異なる。
PJM、CAISO、ERCOT などの ISO/RTO が市場運営を担い、
需給調整市場や容量市場が高度に発達している。

特にテキサス州（ERCOT）は完全自由化市場として知られ、
2021 年の寒波時にはスポット価格が 9,000 \$/MWh に達するなど、
極端な価格スパイクが発生した。

% --- アジアモデル ---
\subsubsection{（3）アジアモデル：段階的自由化と国家主導}

日本・韓国・台湾などが段階的自由化を進めている。
再エネ導入は政策主導で急速に進む一方、市場設計は欧米に比べて未成熟であり、
制度変更リスクが大きい。

日本では 2016 年の小売全面自由化以降、新電力のシェアが約 20\% に達したが、
2021 年の価格高騰時には多数の新電力が経営危機に直面した。

% --- 新興国モデル ---
\subsubsection{（4）新興国モデル：国家主導型市場化}

中国・インド・ブラジルなどでは国家主導で市場化が進む。
再エネ導入は急速だが、価格規制や国家主導の投資が中心であり、
自由化の度合いは限定的である。

% --- 四モデルの比較 ---
\subsubsection{（5）四つのモデルの比較と新電力リスクへの含意}

欧州：価格変動リスクが大きい  
北米：地域差が極端  
アジア：制度リスクが大きい  
新興国：政策依存度が高い  

日本はアジアモデルに属し、段階的自由化と国家主導の制度設計のもとで市場が形成されているが、
再エネ比率の低さや市場設計の未成熟さから、外生ショックに対して脆弱な構造を持つ。

% ============================================================
% 1.2 日本の電力システム改革と市場構造（冒頭）
% ============================================================

\section{日本の電力システム改革と市場構造}

日本の電力システム改革は、1995 年の電気事業法改正に端を発し、
発電・送配電・小売の三部門を段階的に自由化する形で進展してきた。
本節では、1995 年以前の地域独占体制から、2016 年の小売全面自由化、
2020 年の送配電分離に至るまでの制度改革を体系的に整理する。

% ============================================================
% 1.2 日本の電力システム改革と市場構造（拡張版・続き）
% ============================================================

\subsection{1995 年以前の日本の電力システム：地域独占と総括原価方式}

1951 年の電気事業再編により、北海道から沖縄までの 10 電力会社が
地域独占の形で発電・送電・配電・小売を一体的に担う体制が確立した。
この体制は、戦後復興期から高度経済成長期にかけて、
電力の安定供給と大規模投資の実行に一定の成果を上げた。

しかし、1980 年代後半以降、以下の構造的課題が顕在化した。

\begin{itemize}
  \item \textbf{競争の欠如}：地域独占により効率化インセンティブが弱い。
  \item \textbf{料金の硬直性}：総括原価方式によりコストがそのまま料金に転嫁される。
  \item \textbf{電源構成の硬直化}：原子力・大規模火力中心で再エネ導入が進まない。
  \item \textbf{需給調整の非効率性}：広域連系線の容量が小さく地域間融通が限定的。
\end{itemize}

これらの課題は、国際的な自由化潮流と国内の電気料金高騰を背景に、
制度改革の必要性を高める要因となった。

% ------------------------------------------------------------
% 1.2.1 世界的潮流が日本に与えた影響
% ------------------------------------------------------------

\subsection{世界的潮流が日本に与えた影響：自由化・脱炭素・市場化の圧力}

1990 年代以降、欧米を中心に電力市場の自由化が急速に進展した。
英国では 1990 年の電力法改正により発電・送電・小売が分離され、
米国でも ISO/RTO による市場運営が一般化した。
EU は 1996 年の電力指令を皮切りに域内市場の統合と完全自由化を段階的に実施した。

こうした国際的潮流は日本にも強い圧力を与えた。

\begin{itemize}
  \item \textbf{国際比較で高水準の電気料金}  
        日本の産業用電気料金は OECD 諸国の中でも高く、製造業の競争力低下が懸念された。
  \item \textbf{再エネ導入の遅れ}  
        欧州では風力・太陽光が急速に普及した一方、日本では制度的制約により普及が進まなかった。
  \item \textbf{市場の透明性不足}  
        世界銀行や OECD は日本の地域独占・総括原価方式を問題視し、競争導入を提言した。
\end{itemize}

国内でも産業界（経団連）が電気料金引き下げと制度改革を強く要求し、
政府は制度改革を避けられない状況に置かれた。

% ------------------------------------------------------------
% 1.2.2 政策形成プロセス
% ------------------------------------------------------------

\subsection{電力システム改革の政策形成：経産省・電力会社・産業界・有識者会議}

日本の電力制度改革は、政府・電力会社・産業界・有識者会議が相互に影響を及ぼす
政策過程の中で形成された。

\begin{itemize}
  \item \textbf{通商産業省（現・経済産業省）}  
        「電気事業制度改革研究会」「電気事業審議会」を設置し、制度改革の方向性を検討。
  \item \textbf{電力会社（一般電気事業者）}  
        地域独占維持を主張しつつも、国際潮流を受け段階的自由化を受け入れた。
  \item \textbf{産業界（経団連）}  
        電気料金の国際比較に基づき自由化を強く要求。
  \item \textbf{有識者会議}  
        市場透明性向上、競争導入、広域需給調整の必要性を提言。
\end{itemize}

これらの議論を踏まえ、1995 年の電気事業法改正が実施された。

% ------------------------------------------------------------
% 1.2.3 三段階自由化
% ------------------------------------------------------------

\subsection{改正電気事業法と三段階自由化：発電・小売・送配電の再設計}

1995 年以降、日本の電力制度改革は以下の三段階で進展した。

\subsubsection{（1）第一段階：1995 年改正電気事業法による部分自由化}

特定規模需要（大口需要家）に対する部分的な小売自由化が導入され、
発電事業は届出制となり、IPP が参入した。

\subsubsection{（2）第二段階：2000 年代の発電自由化と小売自由化の拡大}

2000 年に卸電力取引所（JEPX）が設立され、発電部門の競争が市場メカニズムを通じて進展。
2004〜2005 年の改正で小売自由化の対象が拡大し、高圧需要家の大部分が自由化市場へ移行した。

\subsubsection{（3）第三段階：2016 年の小売全面自由化と 2020 年の送配電分離}

2016 年に家庭用を含むすべての需要家が自由化市場で電力会社を選択可能となり、
多数の新電力が参入した。

2020 年には送配電部門の法的分離が実施され、ネットワークの中立性が制度的に担保された。

% ------------------------------------------------------------
% 1.2.4 自由化後の市場構造の変容
% ------------------------------------------------------------

\subsection{自由化後の市場構造の変容：競争導入・価格形成・市場成熟度}

\subsubsection{（1）競争導入：新電力の参入と市場シェアの変化}

小売電気事業者の登録数は 2020 年に 700 社を超え、
新電力の販売電力量シェアは 2023 年に約 20\% に達した。

都市部ではスイッチング率が高く、東京電力管内では家庭用の約 25\% が新電力へ移行した。

\subsubsection{（2）価格形成：JEPX を中心とした市場価格の変動性}

JEPX のスポット価格は、再エネ比率の上昇、燃料価格の変動、需給逼迫に応じて大きく変動する。
2021 年 1 月にはスポット価格が 200 円/kWh を超える異常値を記録し、
多数の新電力が調達コスト急騰に直面した。

\subsubsection{（3）市場成熟度：競争の実効性と制度的課題}

旧一般電気事業者の市場支配力、託送料金制度の複雑さ、再エネ接続制約、
家庭用需要家の価格感応度の低さなど、競争市場としての成熟度には課題が残る。

% ============================================================
% 1.3 新電力のビジネスモデル（完全版）
% ============================================================

\section{新電力のビジネスモデル}

本節では、新電力（小売電気事業者）のビジネスモデルを、調達・需給管理・販売・回収という4つの主要機能に分解し、その構造的特徴を整理する。

% ---------------------------------------------
% 図1（ランドスケープ＋TikZ）
% ---------------------------------------------
\begin{landscape}
\begin{figure}[H]
\centering
\begin{tikzpicture}[
    box/.style={
        rectangle,
        draw,
        rounded corners,
        align=center,
        minimum width=48mm,
        minimum height=12mm,
        font=\small
    },
    smallbox/.style={
        rectangle,
        draw,
        rounded corners,
        align=center,
        minimum width=40mm,
        minimum height=12mm,
        font=\small
    },
    arrow/.style={->, thick}
]

% --- 上段（中央に揃える） ---
\node[box] (policy) at (6,7) {制度設計・規制\\（経産省・OCCTO）};
\node[box] (cost)   at (6,5) {制度コスト\\（需給調整金・容量市場）};

% --- 右上（外部ショック） ---
\node[box] (shock)  at (13,6.5) {外部ショック\\（燃料価格・気象・需給逼迫）};

% --- 左側 ---
\node[box] (gen1)   at (-2,3) {大規模発電事業者\\（例：東京電力・JERA等）};
\node[box] (jepx)   at (-2,1) {卸電力市場（JEPX）\\（スポット・先渡・長期契約）};
\node[box] (fit)    at (-2,-1) {FIT電源\\（固定価格買取制度）};
\node[box] (selfgen)at (-2,-3) {自社電源\\（太陽光・風力・バイオマス等）};

% --- 中央 ---
\node[box] (retailer) at (6,1) {新電力（小売電気事業者）\\（需給管理・インバランス管理）};

% --- 右側（横幅を縮めた顧客ボックス） ---
\node[smallbox] (customer) at (14,1)
{顧客\\（料金プラン・地域サービス・\\再エネメニュー）};

% --- 下段 ---
\node[box] (imbalance) at (6,-2.5) {インバランス料金\\（需給誤差ペナルティ）};
\node[box] (risk)      at (6,-5) {新電力の損益・倒産リスク};

% ============================================================
%  矢印（あなたの指定どおりに完全再現）
% ============================================================

% --- 左側：大規模発電 → JEPX（下向き） ---
\draw[arrow] (gen1.south) -- (jepx.north) node[midway, right] {供給};

% --- 左側：FIT → JEPX（上向き） ---
\draw[arrow] (fit.north) -- (jepx.south) node[midway, right] {市場へ放出};

% --- 左側：自社電源 → 新電力（斜め） ---
\draw[arrow] (selfgen.east) -- (retailer.south west);

% --- 左側：JEPX → 新電力（右向き） ---
\draw[arrow] (jepx.east) -- (retailer.west) node[midway, above] {外部調達};

% --- 上段：制度設計 → 制度コスト（下向き） ---
\draw[arrow] (policy.south) -- (cost.north);

% --- 上段：制度コスト → 新電力（下向き、距離短く） ---
\draw[arrow] (cost.south) -- (retailer.north) node[midway, right] {負担};

% --- 外部ショック → 新電力（斜め） ---
\draw[arrow] (shock.south west) -- (retailer.north east);

% --- 中央：新電力 → インバランス料金（下向き） ---
\draw[arrow] (retailer.south) -- (imbalance.north) node[midway, right] {需給誤差};

% --- 下段：インバランス料金 → 損益・倒産リスク（下向き） ---
\draw[arrow] (imbalance.south) -- (risk.north);

% --- 新電力 ↔ 顧客（双方向） ---
\draw[arrow] (retailer.east) -- (customer.west) node[midway, above] {販売};
\draw[arrow] (customer.west) -- (retailer.east) node[midway, below] {料金回収};

\end{tikzpicture}

\caption{新電力のビジネスモデルの全体構造図}
\label{fig:business_model}
\end{figure}
\end{landscape}

% ------------------------------------------------------------
% 1.3.1 三層モデル
% ------------------------------------------------------------

\subsection{新電力ビジネスの全体構造：調達・需給管理・販売の三層モデル}

新電力のビジネスモデルは、図 \ref{fig:business_model}に示すように、

\begin{itemize}
  \item （1）電力の調達  
  \item （2）需給管理  
  \item （3）販売（小売）
\end{itemize}

の三つの主要プロセスから構成される。

以下、各層を詳細に検討する。

% ------------------------------------------------------------
% 1.3.2 調達
% ------------------------------------------------------------

\subsection{調達：発電事業者・卸電力市場・FIT 電源の位置づけ}

調達は新電力の収益構造を決定する最重要基盤である。

\subsubsection{電気という商品の同質性と託送の必然性}

電気は完全に同質的な商品であり、差別化は困難である。
そのため新電力は「電気以外の価値」で競争せざるを得ない。

\subsubsection{（1）発電事業者との相対契約}

価格の固定性・供給量の確実性が高いが、柔軟性に欠ける。

\subsubsection{（2）卸電力市場（JEPX）での取引}

スポット・先渡・時間前市場を活用し、不足分を補う。
価格変動リスクが大きく、2021 年の高騰時には多数の新電力が経営危機に直面した。

\subsubsection{（3）FIT 電源の位置づけ}

価格は固定だが供給量が不確実であり、需給管理の難易度を高める。

\subsubsection{（4）調達ポートフォリオとリスク管理}

価格リスク・量的不確実性・インバランスリスクを最小化するため、
相対契約・市場調達・FIT を組み合わせる。

% ------------------------------------------------------------
% 1.3.3 需給管理
% ------------------------------------------------------------

\subsection{需給管理：インバランス制度と広域機関の役割}

需給管理は新電力の中核機能であり、事業の安定性を左右する。

\subsubsection{（1）計画値同時同量と需要予測}

需要予測の精度が低いとインバランスが増加し、経営リスクが高まる。

\subsubsection{（2）インバランス制度}

不足インバランスは高額であり、2021 年には 200 円/kWh を超える事例も発生した。

\subsubsection{（3）広域機関（OCCTO）の役割}

需給計画の集約、広域調整、インバランス精算ルールの策定などを担う。

\subsubsection{（4）一般送配電事業者の役割}

スマートメーターによる実績値計測、インバランス精算、周波数維持など。

\subsubsection{（5）需給管理の経営的意義}

インバランスリスク・調達リスク・信用リスクを左右する。

% ------------------------------------------------------------
% 1.3.4 販売
% ------------------------------------------------------------

\subsection{販売：料金メニュー・付加価値サービス・顧客獲得戦略}

販売は収益を生み出す最終段階である。

\subsubsection{（1）料金メニューの多様化}

従量料金、時間帯別料金、セット割、再エネメニューなど。

\subsubsection{（2）付加価値サービス}

ポイント還元、スマートホーム連携、環境価値サービスなど。

\subsubsection{（3）顧客獲得戦略}

価格比較サイト、セット割、環境価値訴求、地域密着型サービスなど。

\subsubsection{（4）販売部門の収益構造とリスク}

調達コスト急騰、顧客離脱、信用リスクなど。

\subsubsection{（5）販売戦略の総括}

販売戦略は調達・需給管理と密接に連動し、柔軟な見直しが必要。

% ============================================================
% 1.4 卸電力市場の詳細（完全版）
% ============================================================

\section{卸電力市場の詳細}

本節では、日本の卸電力市場（JEPX）を中心に、市場規模、価格推移、参加主体、
制度的要因などを定量データに基づいて整理する。
1.3 で示した新電力のビジネスモデルは、卸電力市場の価格変動や需給状況に強く依存しており、
市場の実態を把握することは新電力のリスク構造を理解する上で不可欠である。

\subsection{市場規模と取引量の推移（2020--2024）}

JEPX の年間取引量（TWh）、スポット市場のシェア、先渡市場の取引量などを示し、
自由化以降の市場拡大の実態を明らかにする。
また、取引量の季節性や再エネ出力との関連についても整理する。

\begin{figure}[H]
\centering
\includegraphics[width=14cm]{../figure/jepx_volume_2020_2024.pdf}
\caption{JEPX 年間取引量の推移（2020--2024）\\
{\small 出典：JEPX 公表データ（年間取引量統計）}}
\label{fig:jepx_volume}
\end{figure}



\subsection{市場参加者数の推移（新電力・発電事業者）}

小売電気事業者（新電力）、発電事業者、卸売事業者の登録数の推移を示し、
市場参入の動向と競争環境の変化を整理する。

\begin{figure}[H]
\centering
\includegraphics[width=14cm]{../figure/jepx_participants_2020_2024.pdf}
\caption{JEPX 市場参加者数の推移（2020--2024）\\
{\small 出典：JEPX 公表データ（市場参加者数統計）}}
\label{fig:jepx_participants}
\end{figure}

\subsection{退出企業数の推移（2020--2024）}

図 \ref{fig:exit_count} は、小売電気事業者の退出企業数の推移を示したものである。
2021 年以降、卸電力市場価格の高騰やインバランス料金の上昇を背景に、
退出企業数が急増していることが分かる。
退出には、法的倒産、実質的倒産、事業撤退、事業譲渡など複数の形態が含まれるが、
これらの分類および定義については次節（1.5 節）で詳述する。

\begin{figure}[H]
\centering
\includegraphics[width=14cm]{../figure/exit_count_2020_2024.pdf}
\caption{小売電気事業者の退出企業数の推移（2020--2024）\\
{\small 出典：経済産業省・電力広域的運営推進機関（OCCTO）公表資料を基に作成}}
\label{fig:exit_count}
\end{figure}


\subsection{スポット価格の推移と変動性}

スポット市場の平均価格、中央値、価格分布（ヒストグラム）、
および 2021 年 1 月の高騰事例を含む時系列推移を示す。
燃料価格、需給逼迫、再エネ出力の変動との関連も分析する。

\begin{figure}[H]
\centering
\includegraphics[width=15cm]{../figure/jepx_spot_timeseries_2020_2024.pdf}
\caption{JEPX スポット市場価格の推移（2020--2024）\\
{\small 出典：JEPX 公表データ（スポット市場価格統計）}}
\label{fig:jepx_spot_timeseries}
\end{figure}

\begin{figure}[H]
\centering
\includegraphics[width=14cm]{../figure/jepx_spot_hist_2020_2024.pdf}
\caption{JEPX スポット市場価格の分布（2020--2024）\\
{\small 出典：JEPX 公表データ（スポット市場価格統計）}}
\label{fig:jepx_spot_hist}
\end{figure}


\subsection{インバランス料金の推移}

インバランス料金の制度変更（2021 年 4 月）前後の比較、
高騰時のピーク値、季節性などを示し、
新電力の経営リスクとの関連を整理する。

\begin{figure}[H]
\centering
\includegraphics[width=15cm]{../figure/imbalance_timeseries_2020_2024.pdf}
\caption{インバランス料金の推移（2020--2024）\\
{\small 出典：OCCTO・電力広域的運営推進機関 公表データ}}
\label{fig:imbalance_timeseries}
\end{figure}


\subsection{需給逼迫アラート発令回数の推移}

2021 年以降の需給逼迫アラート（旧：需給ひっ迫注意報／警報）の発令回数を整理し、
需給逼迫が市場価格に与える影響を定量的に示す。
また、アラート発令日とスポット価格の関係も分析する。

\begin{figure}[H]
\centering
\includegraphics[width=14cm]{../figure/alert_count_2020_2024.pdf}
\caption{需給逼迫アラート発令回数の推移（2020--2024）\\
{\small 出典：経済産業省・電力広域的運営推進機関（OCCTO）公表資料}}
\label{fig:alert_count}
\end{figure}

\begin{figure}[H]
\centering
\includegraphics[width=14cm]{../figure/alert_vs_price_scatter.pdf}
\caption{アラート発令日とスポット価格の関係\\
{\small 出典：JEPX・OCCTO 公表データを基に作成}}
\label{fig:alert_vs_price_scatter}
\end{figure}


\subsection{JEPX の市場構造（スポット・先渡・時間前）}

各市場の役割、取引量の割合、価格形成への寄与を整理する。
また、先渡市場の流動性の低さや時間前市場の重要性についても説明する。

\begin{figure}[H]
\centering
\begin{tikzpicture}[node distance=2.2cm, >=stealth, thick]

% Nodes
\node[draw, rounded corners, fill=blue!10, minimum width=4cm, minimum height=1cm] (spot) {スポット市場};
\node[draw, rounded corners, fill=green!10, minimum width=4cm, minimum height=1cm, left=of spot] (fwd) {先渡市場（低流動性）};
\node[draw, rounded corners, fill=orange!10, minimum width=4cm, minimum height=1cm, right=of spot] (intraday) {時間前市場（調整）};

\node[draw, rounded corners, fill=gray!10, minimum width=4cm, minimum height=1cm, above=1.8cm of spot] (gen) {発電事業者};
\node[draw, rounded corners, fill=gray!10, minimum width=4cm, minimum height=1cm, below=1.8cm of spot] (retail) {小売電気事業者（新電力）};

% Arrows from generator
\draw[->] (gen) -- node[right]{供給計画} (spot);
\draw[->] (gen) -- node[above]{ヘッジ契約} (fwd);
\draw[->] (gen) -- node[above]{調整} (intraday);

% Arrows to retailer
\draw[->] (spot) -- node[right]{実需に応じた調達} (retail);
\draw[->] (fwd) -- node[above]{価格変動リスクの低減} (retail);
\draw[->] (intraday) -- node[above]{需給調整} (retail);

% Notes
\node[align=left, text width=6cm, below=0.8cm of retail] (note1)
{\small ・先渡市場は流動性が低く、十分なヘッジが困難\\
 ・時間前市場は需給調整の要だが、価格変動が大きい\\
 ・スポット市場への依存が高いほど、価格ショックに脆弱};

\end{tikzpicture}

\caption{JEPX の市場構造（スポット・先渡・時間前）\\
{\small 出典：電力広域的運営推進機関（OCCTO）資料を基に作成}}
\label{fig:jepx_market_structure}
\end{figure}

\subsection{価格形成メカニズム}

スポット市場の需給曲線、発電コスト構造、燃料価格、再エネ出力、
FIT 電源の市場放出量など、価格形成に影響を与える要因を体系的に整理する。

\begin{figure}[H]
\centering
\begin{tikzpicture}[scale=1.1, >=stealth, thick]

% Axes
\draw[->] (0,0) -- (7,0) node[right]{需要量};
\draw[->] (0,0) -- (0,6) node[above]{価格};

% Demand curve
\draw[blue, thick] (1,5) -- (6,1) node[midway, above right]{需要曲線};

% Supply curve (normal)
\draw[red, thick] (1,1) -- (6,5) node[midway, above left]{供給曲線};

% Supply curve shifted (renewables)
\draw[red!50, dashed, thick] (0.5,1) -- (5.5,5)
    node[midway, above left]{再エネ増加による供給曲線の左シフト};

% Intersection (equilibrium)
\filldraw[black] (3.5,3) circle (2pt);
\node[right] at (3.5,3) {均衡価格（スポット価格）};

% Cost blocks (baseload, mid, peak)
\draw[fill=green!20, opacity=0.6] (0.2,0.2) rectangle (2.5,0.8);
\node at (1.3,0.5) {\small ベースロード};

\draw[fill=yellow!20, opacity=0.6] (2.5,0.2) rectangle (4.5,1.2);
\node at (3.5,0.8) {\small ミドル電源};

\draw[fill=red!20, opacity=0.6] (4.5,0.2) rectangle (6.5,2.0);
\node at (5.5,1.4) {\small ピーク電源};

% FIT supply block
\draw[fill=blue!20, opacity=0.4] (0.2,1.2) rectangle (1.5,1.7);
\node at (0.85,1.45) {\tiny FIT電源};

\end{tikzpicture}

\caption{需給曲線の概念図（価格形成メカニズム）\\
{\small 出典：電力広域的運営推進機関（OCCTO）資料を基に作成}}
\label{fig:demand_supply_curve}
\end{figure}



\subsection{市場参加主体の役割}

電力市場には、発電事業者、小売電気事業者、卸売事業者、アグリゲーターなど、多様な主体が参加しており、それぞれが異なる目的と制約のもとで行動している。発電事業者は、燃料価格や設備稼働状況に応じて供給計画を策定し、市場に電力を供出する主体である。一方、小売電気事業者は、需要家との契約に基づき必要量を調達し、価格変動リスクを管理しながら販売を行う。卸売事業者は、需給調整やポートフォリオ最適化を目的として市場間で電力を仲介し、流動性の確保に寄与する。また、アグリゲーターは分散型電源や需要側リソースを束ね、需給調整市場や容量市場に参加することで、系統安定化に重要な役割を果たす。

これらの主体の行動は、スポット市場価格や需給逼迫リスクに直接的な影響を与える。特に、新電力は資本力や調達手段の制約からスポット市場への依存度が高く、他主体の行動変化が経営リスクに直結しやすい構造を持つ。したがって、市場参加主体の役割と相互作用を理解することは、価格形成メカニズムや倒産リスクの分析に不可欠である。

\subsection{制度的要因と市場への影響}

電力市場の価格形成や調達コストには、FIT 電源の市場放出、容量市場、需給調整市場など、制度的要因が大きく影響している。FIT 電源は固定価格買取制度のもとで導入が進んだ再生可能エネルギーであり、その市場放出量は供給曲線を左シフトさせ、スポット価格の低下圧力として作用する。一方、容量市場は将来の供給力確保を目的とした制度であり、発電事業者に対する固定的な収入源となるが、小売電気事業者にとっては追加的なコスト負担となる。

さらに、需給調整市場は系統安定化のための調整力を確保する仕組みであり、調整力の不足や制度変更に伴うコスト上昇は、小売電気事業者の調達コストを押し上げる要因となる。特に新電力は、調整力確保のためのリソースが限られているため、制度的コストの増加が経営に与える影響が相対的に大きい。

このように、制度的要因は単に市場価格に影響を与えるだけでなく、事業者間の競争構造やリスク分担のあり方を規定する重要な要素である。制度設計の変化が新電力の経営環境にどのような影響を及ぼすかを把握することは、倒産リスクの構造的理解に不可欠である。
。

以上のように、本節では日本の卸電力市場の制度的枠組みと価格形成の特徴を整理した。これらの市場構造は、新電力会社が直面する外部環境を規定するものであり、次節で扱う「新電力のリスク構造」を理解するための基盤となる。


% ============================================================
% 1.5 倒産・退出・継続の定義（改訂版）
% ============================================================

\section{倒産・退出・継続の定義}

新電力企業の倒産リスクを分析するためには、倒産・退出・継続という概念を
明確に区別する必要がある。本節では、倒産研究の主要文献
（Beaver, Altman, Ohlson, Shumway）を踏まえつつ、
日本の倒産法制および電力市場の制度的特徴を反映した定義を整理する。

\subsection{法的倒産}

破産、民事再生、会社更生など、裁判所による法的手続きが開始された状態を指す。
最も厳密かつ客観的に観測可能な倒産形態であり、官報や裁判所記録に基づき確認できる。

\subsection{実質的倒産}

法的手続きには至らないものの、企業が事業継続不能な状態に陥った場合を指す。
具体例として、供給停止、債務不履行、小売電気事業の廃止届提出などが含まれる。
新電力市場では、外部ショックにより短期間で資金繰りが悪化し、
この実質的倒産に至る事例が多い。

\subsection{市場退出}

倒産に限らず、事業撤退、事業譲渡、M\&A による吸収など、
企業が市場から退出する一連の事象を指す。
市場退出は必ずしも財務破綻を伴うわけではなく、
戦略的撤退や事業再編の結果として生じる場合もある。

\subsection{事業継続}

供給義務を果たしつつ、財務的に継続可能な状態を指す。
新電力企業の場合、需給管理の安定性、調達戦略、資金調達力などが
事業継続性を左右する。

\subsection{事業単位と企業単位の区別}

新電力は「事業単位」で参入・退出するが、倒産は「企業単位」で発生する。
例えば、事業撤退（廃止届）は市場退出に該当するが、
企業全体が倒産したとは限らない点に注意が必要である。

\subsection{供給停止}

実質的倒産の代表的形態であり、需給管理の失敗や調達費用の急騰により発生する。
新電力市場では、外部ショックにより供給停止が倒産の直接要因となる事例が多い。

\subsection{廃止届}

小売電気事業の廃止届は、事業撤退の公式記録であるが、
財務破綻を伴わない場合もある。倒産と市場退出を区別するための重要な情報源となる。

\vspace{1\baselineskip}
以上の定義は個別には理解しやすいものの、倒産・退出・継続の関係性は
しばしば混同されやすい。そこで図 \ref{fig:bankruptcy_structure} では、
これらの概念がどのように階層的に整理され、互いにどのような位置づけを
持つのかを視覚的に示している。特に、新電力市場では法的倒産と実質的倒産、
事業撤退が同時に観察されるため、概念構造を明確に把握することが重要である。

\begin{figure}[H]
\centering
\begin{tikzpicture}[
  scale=0.78,
  transform shape,
  node distance=12mm,
  box/.style={rectangle, draw, rounded corners, align=center, minimum width=32mm, minimum height=10mm}
]

% Top level
\node[box] (universe) {市場退出（Exit）};

% Second level
\node[box, below left=of universe] (legal) {法的倒産\\(破産・民事再生・会社更生)};
\node[box, below=of universe] (failure) {実質的倒産\\(供給停止・債務不履行)};
\node[box, below right=of universe] (withdraw) {事業撤退\\(廃止届・事業譲渡)};

% Third level (legal breakdown)
\node[box, below left=of legal] (hasan) {破産};
\node[box, below=of legal] (saisei) {民事再生};
\node[box, below right=of legal] (kosei) {会社更生};

% Going concern
\node[box, right=35mm of universe] (going) {事業継続\\(Going Concern)};

% Arrows
\draw[->] (universe) -- (legal);
\draw[->] (universe) -- (failure);
\draw[->] (universe) -- (withdraw);

\draw[->] (legal) -- (hasan);
\draw[->] (legal) -- (saisei);
\draw[->] (legal) -- (kosei);

\end{tikzpicture}
\caption{倒産・退出・継続の概念構造}
\label{fig:bankruptcy_structure}
\end{figure}


% ============================================================
% 1.6 新電力のリスク構造（改訂版）
% ============================================================

\section{新電力のリスク構造}

新電力企業の倒産リスクは、単一の要因ではなく、
市場環境・制度設計・事業運営・財務構造・顧客行動といった
複数の要因が相互に作用することで形成される。
本節では、新電力が直面する主要リスクを体系的に整理し、
倒産リスク分析の基礎枠組みを提示する。

\subsection{市場リスク}

JEPX 価格の変動、需給逼迫、燃料価格の急騰など、
外部環境の変化が企業の調達コストに直接波及するリスクである。
新電力はアセットライト型であるため、市場変動の影響を受けやすい。

\subsection{制度リスク}

インバランス制度、容量市場、FIT 制度、需給調整金など、
制度変更や制度コストの増加が企業キャッシュフローを圧迫するリスクである。
制度コストの急増は倒産ラインを押し上げる要因となる。

\subsection{オペレーショナルリスク}

需給管理の誤差、システム障害、業務運営の不備など、
事業運営上のミスや不確実性に起因するリスクである。
需給管理の精度は新電力の経営安定性に直結する。

\subsection{財務リスク}

資金繰り、負債構成、親会社支援、金融機関からの借入など、
企業の財務構造に起因するリスクである。特に金融環境が悪化した局面では、
銀行の与信姿勢が急速に厳格化し、信用収縮（credit crunch）、
貸し渋り（credit tightening）、貸し剥がし（credit withdrawal）といった
現象が発生する。これらは企業の資金調達力を直接的に低下させ、
外部ショック時の倒産リスクを大きく押し上げる要因となる。

\subsection{顧客リスク}

顧客の解約、未収金、需要変動など、
販売部門に起因するリスクである。
顧客離脱は収益の急減につながり、財務リスクを高める。

\vspace{1\baselineskip}
新電力企業の倒産リスクは、複数のリスク要因が相互に作用しながら形成される。
図 \ref{fig:risk_network} は、市場リスク・制度リスク・オペレーショナルリスク・
顧客リスク・財務リスクの五つの主要リスクがどのように倒産リスクへ影響し、
また財務リスクを媒介として間接的に増幅されるのかを整理したものである。
本図により、倒産リスクが単一要因ではなく、構造的なネットワークとして
理解されるべきであることが明確になる。

\begin{figure}[htbp]
  \centering
  \begin{tikzpicture}[
    node distance=18mm,
    risk/.style={rectangle, draw, rounded corners, align=center, minimum width=24mm, minimum height=9mm},
    arrow/.style={->, thick},
    darrow/.style={->, thick, dashed}
  ]

    % 上段：4つのリスク
    \node[risk] (market) {市場リスク};
    \node[risk, right=of market] (institution) {制度リスク};
    \node[risk, right=of institution] (operational) {オペレーショナル\\リスク};
    \node[risk, right=of operational] (customer) {顧客リスク};

    % 中央下段：財務リスク
    \node[risk, below=18mm of institution] (financial) {財務リスク};

    % 右端：倒産リスク
    \node[risk, right=32mm of financial] (default) {倒産リスク};

    % 各リスクから倒産リスクへの矢印（直接の影響）
    \draw[arrow] (market.east) to[bend left=10] (default.north west);
    \draw[arrow] (institution.south east) to[bend left=5] (default.north);
    \draw[arrow] (operational.south) to[bend left=5] (default.north east);
    \draw[arrow] (customer.east) to[bend left=15] (default.east);
    \draw[arrow] (financial) -- (default);

    % 他リスクから財務リスクへの破線矢印（財務リスクの増幅・媒介）
    \draw[darrow] (market.south) -- (financial.north west);
    \draw[darrow] (institution.south) -- (financial.north);
    \draw[darrow] (operational.south west) -- (financial.north east);
    \draw[darrow] (customer.south west) to[bend right=15] (financial.east);

    % --- 凡例（実線・破線の意味） ---
    \node[align=left, anchor=north east] 
      at ($(current bounding box.south east)+(0,-0.5)$) {
        \begin{tabular}{ll}
        \raisebox{2pt}{\tikz{\draw[->, thick] (0,0) -- (0.8,0);}} 
          & 直接的な影響（倒産リスクに直接作用） \\
        \raisebox{2pt}{\tikz{\draw[->, thick, dashed] (0,0) -- (0.8,0);}} 
          & 財務リスクを通じた間接的な影響（増幅効果） \\
        \end{tabular}
      };

  \end{tikzpicture}
  \caption{新電力における5つのリスクと倒産リスクの関係}
  \label{fig:risk_network}
\end{figure}


% ============================================================
% 1.7 本章のまとめ
% ============================================================

\section{本章のまとめ}

本章では、世界のエネルギー情勢、日本の電力システム改革、新電力のビジネスモデル、
卸電力市場の構造、そして新電力企業が直面するリスク構造を体系的に整理した。

世界的な再エネ拡大、地政学ショック、AI・データセンターによる需要急増といった
外部環境の変化は、自由化市場における価格形成を不安定化させ、
小売事業者の調達リスクを増幅させている。
日本市場では、制度変更や需給調整メカニズムの影響が企業キャッシュフローに
直接的に波及し、外部ショックへの感応度が高い構造が形成されている。
新電力はアセットライト型であるがゆえに、こうした外部環境の変化を
財務構造に吸収しにくい脆弱性を持つ。

また、卸電力市場（JEPX）の価格変動性、インバランス制度の厳格化、
需給調整金の急増といった制度的要因は、新電力企業の倒産ラインを押し上げ、
財務指標の線形的悪化では捉えきれない倒産リスクを形成している。
本章で整理した制度的・市場的背景は、新電力企業が他産業に比べて
外部ショックの影響を強く受ける構造的脆弱性を持つことを示している。

次章では、これらの制度的・市場的背景を踏まえ、
新電力企業の倒産急増という現象を従来モデルでは十分に説明できないという
問題意識から出発し、倒産リスクの構造原理を明確化する。
